%======================================================================
% CAPÍTULO 13 - MODELO DE DATOS
%======================================================================
\chapter{Modelo de Datos}

\bigskip

El modelo de datos de Sentinel define las entidades fundamentales que representan la información de nómina. Este capítulo describe cada entidad y sus relaciones.

\bigskip

\section{Entidades Principales}

El sistema trabaja con tres entidades fundamentales que se fusionan para crear un Expediente Digital Completo (Beneficiario):

\bigskip

\begin{enumerate}
    \item Entidad Base (The Blueprint): Información estructural del afiliado.
    \item Entidad Movimiento (Movements): Estado transaccional dinámico.
    \item Directivas (The Ruleset): Tablas maestras con reglas salariales.
\end{enumerate}

\bigskip

\section{Entidad Beneficiario}

Beneficiario: Entidad central que representa un afiliado militar. Contiene la información de identificación, relaciones con Base y Movimientos, y todos los campos calculados.

\bigskip

La estructura Beneficiario contiene: cedula (identificador único), nombres, apellidos, componente\_id, patterns (clave de fusión), numero\_cuenta, base (datos estructurales), y movimientos (datos financieros).

\bigskip

\section{Entidad Base}

Base: Información estructural del afiliado que contiene los datos invariables: grado, componente, fechas de ingreso y ascenso, antigüedad calculada.

\bigskip

La estructura Base contiene: patterns (clave primaria de búsqueda), grado\_id, componente\_id, fecha\_ingreso, f\_ult\_ascenso, antiguedad (tiempo de servicio calculado en días), antiguedad\_grado (tiempo en el grado actual), st\_no\_ascenso, st\_profesion, sueldo\_base, calculos (primas), garantia\_original, garantia\_anticipo, factor\_aplicado.

\bigskip

\section{Entidad Movimiento}

Movimiento: Estado transaccional dinámico que representa las operaciones financieras del beneficiario: cuentas bancarias, pasivos, y variaciones monetarias.

\bigskip

La estructura Movimiento contiene: cedula (clave de búsqueda), cap\_banco (capital en banco), anticipo, dep\_adicional, dep\_garantia.

\bigskip

\section{Entidades de Directivas}

Las directivas son las reglas salariales que determinan los cálculos:

\bigskip

\begin{center}
\begin{tabular}{llp{8cm}}
\toprule
\textbf{Entidad} & \textbf{Descripción} \\
\midrule
PrimasFunciones & Funciones de prima con fórmulas Rhai, parámetros, condiciones \\
Directiva & Tabla de sueldos base por grado y antigüedad \\
Conceptos & Conceptos de nómina: códigos, descripciones, tipos \\
\bottomrule
\end{tabular}
\end{center}

\bigskip

\section{Relaciones entre Entidades}

Las relaciones entre entidades son las siguientes:

\begin{itemize}
    \item Beneficiario tiene una relación 1:1 con Base.
    \item Beneficiario tiene una relación 1:N con Movimientos.
    \item Base tiene una relación 1:N con Directiva.
\end{itemize}

\bigskip

\section{Flujo de Datos}

El flujo de datos sigue un patrón de procesamiento en pipeline:

\[
\text{Entrada (gRPC)} \rightarrow \text{Deserialización} \rightarrow \text{Cálculo} \rightarrow \text{Fusión} \rightarrow \text{Exportación}
\]

\vfill
\newpage
